\scanpage{43}
[\uline{Ex} $n=1$, $\rho=4-\beta=2-\alpha$, $\sigma=\deleted{\unsure{\alpha}}\;2$, $\widetilde\sigma=1$. Donc

\deleted{Par} 1°) si $\deleted{\unsure{2}}\rho$ inv. \inserted{i.e. $U_\rho\simeq G_m$,} alors $\exists\unsure{1|!}\ u\in U_{\rho=2-\alpha}$ partout $\ne1$, telle que $u$ invarie $f^E$ (ou ce qui revient au $=$, \deleted{$Q_X$} la conique $Q_X$) savoir l'antipodisme, qui est l'unique section $u$ telle que $u^2=1$, $u_s\ne1$ $\forall s\in S$.

2°) Conditions équivalentes

a) $\rho=0$ et \deleted{\unsure{car.|car}} $S$ de car 2 $(2.1_S=0)$

b) $U_\rho$ invarie $f^E$

c) $U_\rho$ \rule{1.5cm}{0.35pt} $Q_X$

d) $\rho=0$ et $\exists$ loc. section partout $\ne1$, qui invarie $f^E$ (ou encore, $Q_X$).]

\begin{center}\rule{4cm}{0.35pt}\end{center}

\uline{Proposition} 1) Si $\sigma\widetilde\rho$ inversible, $\exists1$ section $u$ de $U_\rho$, partout $\ne1$, qui invarie $f^E$ i.e. $Q_X$, savoir l'antipodisme.

2) Conditions équivalentes

a) $\sigma=0$ et $\widetilde\sigma\widetilde\rho=0$ (cette dernière condition conséquence de la première si $n$ pair ou 2 inv)

b) $U_\rho$ invarie $f^{\deleted{\unsure{E|X}}E}$

c) $U_\rho$ invarie $Q_X$

\margin{\unsure{?|2}}d) (\uuline{si $\widetilde\sigma\widetilde\rho=0$}) $\exists$ loc. section de $U_\rho$ partout $\ne1$, qui invarie $f^E$ (\uline{ou} $Q_X$)
% archive sheet number: 42
